\chapter{Getting a Good Bond}
\label{ch:tuning}

This chapter is about what the parameters physically do, and how to move from
a bond that is wrong to one that is right.

\begin{wbwarning}
Change one parameter at a time, in small steps, and bond a few wires after each
change. Changing two at once tells you nothing about which one helped.
\end{wbwarning}

\section{The three that matter most}

Almost all bond quality comes from three things: \textbf{force}, \textbf{power}
and \textbf{energy}.

\begin{description}[leftmargin=!,labelwidth=26mm,style=nextline,font=\normalfont\bfseries]

\item[Force] How hard the tool presses. Too little and the wire moves under the
  tool instead of scrubbing against the pad, so no weld forms. Too much and the
  wire is squashed thin at the heel, which welds well but breaks easily.

\item[Power] How hard the tool scrubs --- the ultrasonic amplitude. This is the
  aggression of the process.

\item[Energy] How much scrubbing there is in total --- the dose. Since the
  machine delivers a fixed energy at a fixed power, raising power shortens the
  weld and raising energy lengthens it.

\end{description}

\begin{wbnote}
Power and energy are separate controls of the same process. A gentle, long weld
(low power, high energy) and an aggressive, short one (high power, low energy)
can deliver the same total energy and produce very different bonds. Low and
long is usually kinder to the pad.
\end{wbnote}

The second bond is made through the wire rather than directly onto it, so it
normally needs more of everything --- which is why the defaults set
\menuitem{Power 2} and \menuitem{Energy 2} above \menuitem{Power 1} and
\menuitem{Energy 1}.

\section{Heights}

\begin{description}[leftmargin=!,labelwidth=26mm,style=nextline,font=\normalfont\ttfamily]

\item[Search 1 / Search 2] Where the fast approach stops and the careful
  descent begins. It only affects cycle time and safety, not bond quality ---
  as long as it is above the workpiece.

\item[Overtravel] How far below the pad the machine aims. It is a push, not a
  destination: the pad stops the tool. It needs to be negative enough that the
  machine keeps pushing after contact, but not so negative that a missed
  touchdown drives the tool hard into the work.

\item[Kink Height] Where the wire is bent as the tool lifts off the first bond.
  Too low and the kink forms in the weld itself, weakening the heel. Too high
  and the loop starts too steeply.

\item[Loop Height] The apex of the loop. Higher gives more clearance and a
  longer wire; lower gives a tighter, flatter loop.

\end{description}

\section{Tail and tear}

\begin{description}[leftmargin=!,labelwidth=26mm,style=nextline,font=\normalfont\ttfamily]

\item[Tail] How much wire is fed after the first bond. It determines how much
  wire protrudes from the tool for the next cycle's first bond.

\item[Tear] How far the axis travels to break the wire after the second bond.
  Too little and the wire does not part; too much and the second bond is
  stressed or lifted.

\end{description}

\begin{wbwarning}
\menuitem{Tail} and \menuitem{Tear} are \emph{relative} travels applied one
after the other, so changing \menuitem{Tail} moves where the tear finishes as
well. After adjusting \menuitem{Tail}, re-check that the wire still tears
cleanly.
\end{wbwarning}

\section{Symptoms and what to change}

Work down each row in order --- the first suggestion is usually the right one.

\begin{center}\footnotesize
\begin{tabular}{@{}L{34mm}L{40mm}L{46mm}@{}}
\toprule
\textbf{Symptom} & \textbf{Try first} & \textbf{Then} \\
\midrule

First bond does not stick &
Raise \menuitem{Energy 1}, then \menuitem{Power 1} &
Raise \menuitem{Force 1 G}. Check the pad is clean and the tool is not worn.
Run \keycap{SETUP} to confirm force is real. \\
\addlinespace

Second bond does not stick &
Raise \menuitem{Energy 2}, then \menuitem{Power 2} &
Raise \menuitem{Force 2 G}. Confirm the loop is not pulling the wire off the
second pad --- try raising \menuitem{Loop Height}. \\
\addlinespace

Bond sticks but lifts in a pull test &
Lower \menuitem{Force} slightly &
Lower \menuitem{Power} and raise \menuitem{Energy} to keep the dose. Check
\menuitem{Kink Height} is not so low that the wire is bent inside the weld. \\
\addlinespace

Wire flattened too thin, breaks at the heel &
Lower \menuitem{Force} &
Lower \menuitem{Energy}. Raise \menuitem{Kink Height}. \\
\addlinespace

Cratering or pad damage &
Lower \menuitem{Power} and raise \menuitem{Energy} to compensate &
Lower \menuitem{Force}. Lower \menuitem{Down Speed} --- impact at touchdown
damages pads independently of the weld. \\
\addlinespace

Tail too short or missing &
Raise \menuitem{Tail} &
Raise \menuitem{Tail Energy}. Check the clamp opens --- watch the
\keycap{CLAMP OPEN} lamp during a cycle. \\
\addlinespace

Tail too long, untidy &
Lower \menuitem{Tail} &
Check \menuitem{Tear} still parts the wire after the change. \\
\addlinespace

Wire does not tear &
Raise \menuitem{Tear} &
Raise \menuitem{Clamp Volt} --- the wire may be slipping in the clamp rather
than failing to tear. \\
\addlinespace

Tearing pulls the second bond off &
Lower \menuitem{Tear} &
Raise \menuitem{Energy 2} so the second bond is stronger than the tear force. \\
\addlinespace

Loop too high or too long &
Lower \menuitem{Loop Height} &
Lower \menuitem{Reverse}. \\
\addlinespace

Loop sagging or shorting &
Raise \menuitem{Loop Height} &
Raise \menuitem{Kink Height} so the loop leaves the first bond more steeply. \\
\addlinespace

Loop shape varies bond to bond &
Lower \menuitem{Up Speed} &
Check the wire spool feeds freely and the clamp voltage is not so high it drags
the wire. \\
\addlinespace

Quality varies with nothing changed &
Run \keycap{TEST}; check frequency and Q &
Run \keycap{SETUP}. Inspect the tool for wear or build-up. \\
\addlinespace

\msg{LOW BONDING POWER} on some bonds &
Widen \menuitem{Scan Start} to \menuitem{Scan Stop} &
Raise \menuitem{Bond Timeout} if welds are genuinely slow; raise
\menuitem{Power} if they are genuinely weak. \\

\bottomrule
\end{tabular}
\end{center}

\section{Developing a new recipe}

\begin{enumerate}[leftmargin=*]
\item Load \menuitem{DEFAULT} and save it under a new name.
\item Set the heights first, in Manual mode if that is easier: get
  \menuitem{Reset Height}, \menuitem{Search}, \menuitem{Overtravel} and
  \menuitem{Loop Height} sensible for the part before touching the weld
  parameters.
\item Run \keycap{TEST} and set the scan window around the reported frequency.
\item Run \keycap{SETUP} so commanded force is real force.
\item Start from the default force, power and energy. Bond a few wires.
\item Tune the first bond until it is consistently good, ignoring the second.
\item Tune the second bond.
\item Tune \menuitem{Tail} and \menuitem{Tear} last --- they depend on
  everything above them.
\item Press \keycap{SAVE}.
\end{enumerate}

\begin{wbnote}
Tune in that order. Tail and tear behaviour depends on how the second bond
came out, so tuning them before the welds are settled means doing it twice.
\end{wbnote}

\section{Starting points by material}

The first three parameter screens follow the 4523AD schedule, and the parameter
names match. \textbf{A bonding schedule developed on a 4523AD --- or published
in the \ksmanual\ --- can be entered on this machine directly.} That is the
fastest way to a working recipe for a material you have not run before.

Two cautions when transferring a schedule:

\begin{center}\small
\begin{tabular}{@{}L{34mm}L{88mm}@{}}
\toprule
\textbf{Parameter} & \textbf{When transferring} \\
\midrule
Force &
  Transfers directly, in grams --- but only after you have run \keycap{SETUP}
  on this machine. Until then the commanded force is uncorrected and a
  transferred figure means little. \\
\addlinespace
Power and energy &
  Transfer as \emph{starting points only}. This machine's ultrasonic drive and
  its transducer are not the ones the schedule was developed on, so the same
  numbers will not produce the same amplitude. Expect to re-tune. \\
\addlinespace
Heights and travels &
  Transfer directly. The mechanics are the same. \\
\addlinespace
Timing, scan, tail assist &
  No equivalent on the 4523AD --- start from the factory defaults in
  Appendix~\ref{app:defaults}. \\
\bottomrule
\end{tabular}
\end{center}

\begin{wbwarning}
Always confirm a transferred schedule on scrap material before production. A
schedule that is right for a 4523AD is a good first guess here, not a
qualified recipe.
\end{wbwarning}

\seeks{published bonding schedules by wire type and diameter, and its bonding
troubleshooting section --- both apply to this machine's mechanics and its
first three parameter screens.}

\TODO{Once recipes are proven here, record them in this section as the
house starting points, superseding the transferred 4523AD values. Give force,
power and energy for both bonds per wire type and pad metallisation.}

\section{Bond quality criteria}

\TODO{Supply the acceptance criteria used here: pull-test force limits per wire
type, acceptable failure modes, sampling frequency, and visual criteria for
bond deformation width and heel condition. Name the standard being worked to if
there is one. The \ksmanual's bond-inspection guidance describes what to look
at; the numeric limits are a local decision and must be stated here.}
